% Part: incompleteness
% Chapter: theories-computability
% Section: first-incompleteness

\documentclass[../../../include/open-logic-section]{subfiles}

\begin{document}

\olfileid{inc}{tcp}{inc}

\olsection{$\Th{Q}$ هېڅ بشپړه، سازګاره او !!^{axiomatizable}
  غځونه نۀ لري}

\begin{thm}
\ollabel{thm:first-incompleteness}
د~$\Th{Q}$ هېڅ بشپړه، سازګاره او !!{axiomatizable} غځونه نشته۔
\end{thm}

\begin{proof}
موږ لا پوهېږو چې د~$\Th{Q}$ هېڅ سازګاره او د پرېکړې وړ غځونه
نشته۔ خو که~$\Th{T}$ بشپړه او !!{axiomatized} وي، نو د پرېکړې وړ
ده۔
\end{proof}

\begin{explain}
دا قضيه د ګوډل د ناتکميلۍ د لومړۍ قضيې له اصلي ۱۹۳۱ز بيان څخه ډېره
لرې نۀ ده۔ له نوې اصطلاح پرته، مهم توپيرونه دا دي: ګوډل د
«سازګار» پر ځاے «$\omega$-سازګار» لري؛ او هغه نۀ شو کولے په بشپړ
عموميت «!!{axiomatizable}» ووايي، ځکه هغه مهال د محاسبه کېدنې صوري
مفهوم لا جوړ شوے نۀ و۔ (د محاسبه کېدنې صوري مدلونه په راتلونکې
لسيزه کښې رامنځته شول، چې ګوډل هم پکې برخه لرله، او تر ډېره د دې
دپاره چې د هغو تيوريو ډولونه ځانګړي کړي چې د ګوډل ښکارندې پرې پلې
کېږي۔)

قضيه وايي چې هر څه يوځای نۀ شئ لرلے؛ يعنې بشپړتيا، سازګاري او
!!{axiomatizability}۔ خو که له دې هر يو پرېږدئ، نور دوه لرلے شئ:
$\Th{Q}$ سازګاره او د محاسبوي بديهي اصولو لرونکې ده، خو بشپړه نۀ
ده؛ ناسازګاره تيوري بشپړه ده او محاسبوي بديهي اصول لري (مثلاً د
$\{ \eq/[0][0] \}$ په وسيله)، خو سازګاره نۀ ده؛ او د حساب د رښتينو
جملو سټ بشپړ او سازګار دے، خو د محاسبوي بديهي اصولو لرونکے نۀ دے۔
\end{explain}
\end{document}
